\begin{figure}[!h]
\centering
\begin{tikzpicture}
\begin{axis}[
    ybar stacked,
    width=0.85\textwidth,
    height=8cm,
    bar width=35pt,
    ylabel={Частка (\%)},
    ymin=0, ymax=105,
    symbolic x coords={Теоретико-аналітичний, Практико-розвивальний, Інтегративно-рефлексивний},
    xtick=data,
    xticklabel style={font=\small, text width=3.5cm, align=center},
    yticklabel style={font=\small},
    legend style={
        at={(0.5,-0.3)},
        anchor=north,
        legend columns=3,
        font=\small,
        draw=none,
    },
    grid=major,
    grid style={gray!20},
    ymajorgrids=true,
    xmajorgrids=false,
]

% Modeling (bottom layer)
\addplot+[ybar, fill=blue!70, draw=blue!90] coordinates {
    (Теоретико-аналітичний, 46.9)
    (Практико-розвивальний, 43.4)
    (Інтегративно-рефлексивний, 38.6)
};
\addlegendentry{Моделювання}

% Scaffolding
\addplot+[ybar, fill=green!60, draw=green!80] coordinates {
    (Теоретико-аналітичний, 4.8)
    (Практико-розвивальний, 7.0)
    (Інтегративно-рефлексивний, 6.7)
};
\addlegendentry{Скафолдинг}

% Questioning
\addplot+[ybar, fill=orange!60, draw=orange!80] coordinates {
    (Теоретико-аналітичний, 14.8)
    (Практико-розвивальний, 11.7)
    (Інтегративно-рефлексивний, 15.7)
};
\addlegendentry{Запитування}

% Feedback
\addplot+[ybar, fill=violet!50, draw=violet!70] coordinates {
    (Теоретико-аналітичний, 16.2)
    (Практико-розвивальний, 15.2)
    (Інтегративно-рефлексивний, 13.4)
};
\addlegendentry{Зворотний зв'язок}

% Troubleshooting
\addplot+[ybar, fill=red!50, draw=red!70] coordinates {
    (Теоретико-аналітичний, 14.8)
    (Практико-розвивальний, 20.6)
    (Інтегративно-рефлексивний, 21.3)
};
\addlegendentry{Вирішення проблем}

% Real-world connection
\addplot+[ybar, fill=cyan!50, draw=cyan!70] coordinates {
    (Теоретико-аналітичний, 2.6)
    (Практико-розвивальний, 2.2)
    (Інтегративно-рефлексивний, 4.3)
};
\addlegendentry{Зв'язок з практикою}

\end{axis}
\end{tikzpicture}
\caption{Розподіл стратегій викладання за етапами методики навчання (\%)}
\label{fig:teaching_strategies}
\end{figure}
